% =====================================================================
%  23.4  Function-Composition Trees.
%  Exports: sec:c29-func  def:c29-func  thm:c29-func-mono  fig:c29-func-auto
%           lst:c29-func  prob:c29-750e  lst:c29-750e
%  The general form of Ch.20's lazy tag: a node's aggregate may be a FUNCTION
%  (affine map, or a min-plus matrix), and the merge is FUNCTION COMPOSITION --
%  which is associative, so the plain Ch.19 fold still works.  Back-cites the
%  Ch.20 composition monoid (\secref{c15-compose}) and DISCHARGES the Ch.21
%  sparse-table boundary forward-cite (\secref{c19-sparse}: "advanced trees ->
%  Ch.23").  The term "function composition" is indexed in this chapter's
%  00_opener, so this file does NOT re-\term it.
%  Solved CF 750E "New Year and Old Subsequence" (2600), live-verified via the
%  problemset API + problem page.  Definition read live: a string is NICE if
%  "2017" occurs as a subsequence but "2016" does NOT; ugliness = min deletions
%  to obtain a nice string, -1 if impossible.  Official samples read live and
%  hand-replayed:
%    S1 "8 3 / 20166766 / q(1,8)(1,7)(2,8)" -> 4,3,-1
%       (greedy automaton on "20166766": delete the four 6's -> 2017; 4)
%    S2 "012016662091670" q(1,14) -> 2 ; q(4,15) -> 1
%    S3 "1234" -> -1,-1
% =====================================================================
\section{Function-Composition Trees}\label{sec:c29-func}

\begin{previouslybox}[Ch.~20's lazy tag was already this idea in disguise]
Chapter~20 taught that a pending tag is a \emph{transform} and that two tags
meeting before a push must \emph{compose} (\(\secref{c15-compose}\)): the
affine pair \((\mathit{mul},\mathit{add})\) merges by \(x\mapsto a_1(a_2x+b_2)+b_1\),
associatively, with an identity.  That was presented as a trick for \emph{deferred
updates}.  Here is the same structure promoted to a full engine: a segment tree
need not store a \emph{number} at each node -- it may store a \emph{function},
and the ``combine the two children'' step is function composition.  Because
composition is associative and has an identity, the entire Chapter~19 build /
merge / query skeleton is unchanged; only the fold operator changes.  When the
functions are affine maps you get range affine-update problems; when they are
\emph{matrices under \((\min,+)\)} you get a whole class of subsequence DP that
no arithmetic fold can express -- CF~750E is the exemplar.
\end{previouslybox}

\begin{topicintro}
A composition tree stores, at the node covering segment \(S\), the single map
\(F_S\) equal to the ordered composition of the per-element maps of \(S\).
Merging is \(F_{[l,r]}=F_{[l,m]}\circ F_{[m+1,r]}\) (order matters -- composition
is not commutative).  Applying the node to a state vector is evaluating \(F_S\)
at that state.  Two families matter:

\emph{Affine maps} \(x\mapsto a\,x+b\), composed as \((a_1,b_1)\circ(a_2,b_2) =
(a_1a_2,\;a_1b_2+b_1)\).  This is exactly Chapter~20's lazy composition monoid,
now used as the \emph{stored aggregate} instead of a deferred tag.

\emph{Min-plus matrices} -- a \(K\times K\) matrix \(M\) acts on a length-\(K\)
state vector \(\mathbf{v}\) by \((\mathbf{v}M)_j=\min_i(v_i+M_{ij})\), and two
matrices compose by \((MN)_{ij}=\min_k(M_{ik}+N_{kj})\).  A matrix is the most
general ``map on a finite set of states that is a \(\min\)/\(+\) combination of
them,'' so any DP whose step is a fixed-size state machine becomes a matrix and
any range of steps becomes a matrix \emph{product} -- answerable by a segment
tree in \(O(K^3\log n)\).

\textbf{Prerequisites.} Chapter~19 (\(\chref{c13}\)) -- the fold skeleton being
reused.  Chapter~20 (\(\secref{c15-compose}\)) -- the composition monoid this
generalises.  Basic matrix multiplication (the \(\min/\!+\) swap is the only
novelty).

\textbf{After this section you can}
\begin{itemize}[nosep]
  \item store an affine or matrix map at a node and merge by composition, keeping
        the \(O(\log n)\) query of Chapter~19; \hfill\lv{4}
  \item explain why associativity of composition (not commutativity) is all the
        fold requires; \hfill\lv{4}
  \item turn a fixed-size subsequence DP into a per-character \((\min,+)\) matrix
        and read a range's answer as a matrix entry; \hfill\lv{5}
  \item solve CF~750E: the ``2017 present, 2016 absent'' automaton as a \(5\times5\)
        min-plus segment tree. \hfill\lv{5}
\end{itemize}
\end{topicintro}

\begin{recognitionbox}[Pattern card -- ``the aggregate is a map, not a number'']
\reckey{a range query whose answer depends on running a \emph{sequential process}
(a DP, an automaton, a recurrence) over the range, where the process has a
\emph{fixed, small} state of size \(K\); merging two adjacent ranges means
``run the left process, then the right'' -- i.e.\ compose their transforms.}
\recuse{
  \item range of affine updates + point query, or range of ``apply \(x\mapsto ax+b\) then add'' -- affine composition tree (\(\defref{c29-func}\));
  \item subsequence / automaton queries over a range (``min edits so that pattern \(\dots\)'') -- \((\min,+)\) matrix automaton tree, e.g.\ CF~750E;
  \item any associative transform, including a linear recurrence step (matrix power over a segment).
}
\recnotuse{
  \item the state grows with the input (unbounded DP) -- no fixed \(K\), no constant-size matrix;
  \item the merge is a simple number fold (\(\sum,\min,\gcd\)) -- plain Ch.~19 tree (\(\chref{c13}\));
  \item the transform is deferred only to push to children and folds a scalar -- that is Ch.~20's lazy tag (\(\secref{c15-compose}\)), same algebra, different role.
}
One-line test: \emph{each element is a fixed-state machine and a range means
chaining the machines -- store the composed map (a matrix) at the node and merge
by \((\min,+)\) product.}
\end{recognitionbox}

% ---------------------------------------------------------------
\subsection{Composition as the fold operator}
% ---------------------------------------------------------------

\begin{definition}{The transform-at-node segment tree}{c29-func}
Let each array element \(i\) carry a map \(f_i\) on a set \(X\), and suppose maps
compose (\((g\circ f)(x)=g(f(x))\)) associatively with an identity
\(\mathit{id}\).  The node covering \([l,r]\) stores \(F_{[l,r]} =
f_r\circ\cdots\circ f_l\), built from children by
\(F_{[l,r]} = F_{[\mathit{mid}+1,r]}\circ F_{[l,\mathit{mid}}\).  A range query
returns the composed map for that range by combining \(O(\log n)\) node maps in
left-to-right order; applying it to an input \(x\in X\) gives the result of
running the whole chain.  For affine maps \(X=\mathbb{R}\), \(f(x)=ax+b\); for a
finite-state machine \(X\) is the set of states and each \(f_i\) is its
transition matrix under \((\min,+)\).
\end{definition}

\begin{theorem}{A composition tree is a segment tree}{c29-func-mono}
If composition \(\circ\) is associative with identity \(\mathit{id}\), the
build/point-update/range-query of \(\defref{c29-func}\) run in \(O(\mathit{cost}(\circ))\)
per level, i.e.\ a range query is \(O(\mathit{cost}(\circ)\log n)\) and the answer
is the exact ordered composition of the range's maps.  No commutativity is
required.
\end{theorem}
\begin{proofbox}
The Chapter~19 skeleton needs only a binary associative combine and a neutral
value for empty nodes -- both supplied by \((\{\text{maps}\},\circ,\mathit{id})\),
a monoid by \(\defref{c29-func}\).  Each internal node stores the composition of
its interval by construction, so a canonical cover of \([l,r]\) into \(O(\log n)\)
maximal nodes, folded \emph{in left-to-right order}, equals \(f_r\circ\cdots\circ
f_l\) by associativity.  Folding in the wrong order would be wrong -- composition
of maps is generally not commutative -- but the tree's canonical-cover traversal
preserves order, so the result is exact. \(\Box\)
\end{proofbox}

\begin{insightbox}[Li Chao, beats, and composition trees are one idea seen thrice]
Chapter~19 folds \emph{numbers} with \(+\) or \(\min\).  A Li Chao tree
(\(\secref{c29-lichao}\)) folds \emph{lines} with ``pointwise best''.  This section
folds \emph{maps} with composition.  All three are the same segment-tree spine
(\(\chref{c13}\)) wearing a different monoid: the axis (positions vs.\
\(x\)-domain) and the combine (sum vs.\ envelope vs.\ compose) vary, but
associativity is the one law that must hold.  Recognising \emph{which monoid} a
problem needs is the skill Chapter~20's \(\secref{c15-compose}\) began and this
chapter completes.
\end{insightbox}

\begin{figure}[H]
\centering
\begin{tikzpicture}[diagram, >=stealth,
  st/.style={circle, draw=primaryblue, fill=primaryblue!12, minimum size=8mm,
             font=\scriptsize\bfseries},
  tr/.style={->, font=\tiny, draw=darkgray}]
  \node[st] (s0) at (0,0)   {\(0\)};
  \node[st] (s1) at (2.3,0) {\(1\)};
  \node[st] (s2) at (4.6,0) {\(2\)};
  \node[st] (s3) at (6.9,0) {\(3\)};
  \node[st, fill=successgreen!18, draw=successgreen] (s4) at (9.2,0) {\(4\)};
  \draw[tr] (s0) to[bend left=20] node[above]{\(\texttt{2}\)} (s1);
  \draw[tr] (s1) to[bend left=20] node[above]{\(\texttt{0}\)} (s2);
  \draw[tr] (s2) to[bend left=20] node[above]{\(\texttt{1}\)} (s3);
  \draw[tr, draw=successgreen] (s3) to[bend left=20] node[above]{\(\texttt{7}\)} (s4);
  \draw[tr, loop below] (s0) node[below, font=\tiny]{other: \(0\)};
  \draw[tr] (s1) to[bend left=20] node[below]{\(\texttt{2}\)} (s0);
  \draw[tr] (s2) to[bend left=20] node[below]{\(\texttt{0},\texttt{2}\)} (s1);
  \draw[tr] (s3) to[bend left=20] node[below]{\(\texttt{1},\texttt{0},\texttt{2}\)} (s2);
  \draw[tr, loop below, draw=warningred] (s3) node[below=1mm, font=\tiny, text=warningred]{\(\texttt{6}\): must delete};
  \draw[tr, loop below, draw=warningred] (s4) node[below=1mm, font=\tiny, text=warningred]{\(\texttt{6}\): must delete};
\end{tikzpicture}
\caption{The CF~750E automaton.  State \(i\) is ``the longest prefix of
\(\texttt{2017}\) that is a subsequence of the kept digits'' has length \(i\).
Reaching state~\(4\) means \(\texttt{2017}\) is present.  A kept digit matching
the next pattern character advances the state (free); any other kept digit is
inert -- \emph{except} a \(\texttt{6}\) at state \(\ge3\), which would create the
forbidden \(\texttt{2016}\) subsequence and so must be deleted (cost~\(1\)).  Each
column of this diagram is one \(5\times5\) min-plus matrix, and a query chains
them.}
\label{fig:c29-func-auto}
\end{figure}

\begin{lstlisting}[caption={Generic min-plus matrix and its composition (the merge operator of a function-composition tree). The identity is the (min,+) identity: 0 on the diagonal, infinity elsewhere. ASCII only}, label={lst:c29-func}]
const int K = 5;                       // number of automaton states
const long long INF = (1LL<<60);
struct Mat {
    long long a[K][K];
    Mat(bool ident=false){
        for (int i=0;i<K;i++) for (int j=0;j<K;j++)
            a[i][j] = (ident ? (i==j?0:INF) : INF);
    }
};
// ordered (min,+) product: apply X first, then Y.  NOT commutative.
Mat compose(const Mat& X, const Mat& Y){
    Mat Z;                              // start all-INF
    for (int i=0;i<K;i++)
        for (int k=0;k<K;k++) if (X.a[i][k] < INF)
            for (int j=0;j<K;j++) if (Y.a[k][j] < INF) {
                long long v = X.a[i][k] + Y.a[k][j];
                if (v < Z.a[i][j]) Z.a[i][j] = v;
            }
    return Z;
}
// affine-map variant (K not needed): (a,b) x |-> a*x+b
// compose (a1,b1) after (a2,b2) = (a1*a2, a1*b2 + b1)
\end{lstlisting}

% ---------------------------------------------------------------
\subsection{Problem Lab: CF 750E -- a subsequence automaton as a matrix}
% ---------------------------------------------------------------

\solvedproblem{c29-750e}{New Year and Old Subsequence}{CF 750E}
  {https://codeforces.com/problemset/problem/750/E}{CF 2600 L5}
  {function-composition tree; min-plus matrix automaton; subsequence-present-and-absent DP as a fixed-state machine; ordered matrix product}

\begin{problemstatement}[CF 750E New Year and Old Subsequence]
A digit string \(s\) of length \(n\) (\(4\le n\le 2\cdot10^5\)).  A string is
\emph{nice} when \(\texttt{2017}\) occurs in it as a subsequence but
\(\texttt{2016}\) does \emph{not}.  The \emph{ugliness} of a string is the fewest
characters that can be deleted to leave a nice string, or \(-1\) if no set of
deletions makes it nice.  Answer \(q\) (\(1\le q\le 2\cdot10^5\)) queries; query
\(i\) gives a range \([a_i,b_i]\) and asks for the ugliness of the substring
\(s[a_i..b_i]\).

\textbf{Input.} Line~1: \(n,q\).  Line~2: \(s\).  Next \(q\) lines: \(a_i\;b_i\).

\textbf{Output.} One ugliness per query.

\textbf{Example (official).}
\begin{verbatim}
Input:
8 3
20166766
1 8
1 7
2 8
Output:
4
3
-1
\end{verbatim}
\end{problemstatement}

\begin{attemptgate}[Attempt it yourself first -- \lv{5}, budget a real hour]
The substring has \(10^5\) choices of deletion set, so you cannot search them.
Ask: what is the \emph{smallest state} that lets me decide, per digit, whether to
keep or delete?  Two subsequences are in play, sharing the prefix \(\texttt{201}\).
\peek{track \(=\) the length of the longest prefix of ``2017'' that is a subsequence of the kept digits so far -- five states \(0..4\).}
\nudge{advancing to state~4 costs nothing (that is the goal), but a kept \(\texttt{6}\) while the state is \(\ge3\) would build the forbidden ``2016'' -- those \(\texttt{6}\)'s are exactly the ones you must delete.}
\fullhint{Per digit \(c\) build a \(5\times5\) min-plus matrix \(M_c\) with
\(M_c[i][i]=1\) (delete \(c\)); if \(c\) equals the next pattern symbol
\(\texttt{2017}[i]\) also set \(M_c[i][i{+}1]=0\) (keep, advance); if \(c=\texttt{6}\)
and \(i\ge3\) leave the keep option \emph{closed}; otherwise set \(M_c[i][i]=0\)
(keep, inert).  A range's answer is the ordered product
\(M_{a}\otimes\cdots\otimes M_{b}\) read at entry \([0][4]\) (\(-1\) if infinite)
-- chain those products with a segment tree over the \(M_c\).}
\end{attemptgate}

\solstep{Understanding}
Nice is a condition on the \emph{kept} subsequence: it must contain
\(\texttt{2017}\) and must avoid \(\texttt{2016}\).  Both strings begin
\(\texttt{201}\); they differ only in the last symbol (\(\texttt{7}\) versus
\(\texttt{6}\)).  So a single running quantity -- the length of the longest
prefix of \(\texttt{2017}\) already a subsequence of the kept digits -- carries
all the information: reaching \(4\) certifies \(\texttt{2017}\), and the state
being \(\ge3\) is exactly when a kept \(\texttt{6}\) would certify the forbidden
\(\texttt{2016}\).  Minimising deletions is then a shortest path through the five
states, and each digit is one \(5\times5\) transition matrix.  The range query is a
matrix product -- a composition tree.

\solstep{Key Observations}
\begin{observation}
Reading a digit \(c\) with current state \(i\): keeping \(c=\texttt{2017}[i]\)
forces \(i\to i{+}1\) at cost \(0\); keeping any other digit leaves \(i\) at cost
\(0\); deleting \(c\) leaves \(i\) at cost \(1\).  Those three rules fill exactly
one column \(M_c\) of \(\figref{c29-func-auto}\).
\end{observation}
\begin{observation}
The forbidden condition turns one keep into a must-delete: while \(i\ge3\) (a
\(\texttt{201}\) subsequence already kept), a \(\texttt{6}\) may only be deleted.
Encode that by setting \(M_c[i][i]=1\) and \emph{not} allowing the \(0\) diagonal
there -- the \((\min,+)\) entry for ``keep and stay'' simply stays infinite.
\end{observation}
\begin{observation}
Concatenating two digit-blocks means applying their matrices in order, and
\((\min,+)\) matrix multiplication is precisely ``best state \(i\) to state
\(j\) through a middle state \(k\)'' -- so the block matrices compose
associatively (\(\defref{c29-func}\)), which is what lets a segment tree store them.
\end{observation}

\solstep{The Approach Ladder}
\begin{center}\renewcommand{\arraystretch}{1.25}\setlength{\tabcolsep}{3pt}
\begin{tabular}{@{}c p{0.26\linewidth} p{0.30\linewidth} l c@{}}
\toprule
\textbf{\#} & \textbf{Approach} & \textbf{Idea} & \textbf{Time} & \textbf{Verdict}\\
\midrule
1 & \lv{3}\ brute force & try every deletion set per query & \(2^{\mathit{len}}\) & \tle \\
2 & \lv{4}\ per-query linear DP & run the 5-state DP over \([a,b]\) & \(\bigO{q\,n\,K}\) & \tle \\
3 & \lv{5}\ \textbf{min-plus matrix segtree} & store/compose \(M_c\) (\(\coderef{c29-750e}\)) & \(\bigO{(n+q)K^3\log n}\) & \ac \\
\bottomrule
\end{tabular}\end{center}
Rung~2 re-walks the substring for every query; the matrix tree does the same DP
\(O(\log n)\) times by storing pre-composed blocks.

\solstep{Solution}
Build, for each position, the \(5\times5\) matrix \(M_{s_i}\) of
Observation~1/2 (state count \(K=5\), pattern \(\texttt{2017}\)).  Load them into
a segment tree whose combine is \(\coderef{c29-func}\)'s ordered \((\min,+)\)
\code{compose}.  For a query \([a,b]\), fold the covering nodes left-to-right into
one matrix \(P\) and output \(P[0][4]\) if finite, else \(-1\).  State~\(0\) is
where every substring starts; state~\(4\) is ``\(\texttt{2017}\) achieved''; the
forbidden \(2016\)'s are already charged as mandatory deletes.

\solstep{Implementation}
\begin{lstlisting}[caption={CF 750E: per-digit min-plus automaton matrices chained by a segment tree. make(c) encodes advance/inert/must-delete-6; query folds the range's matrices in order; answer is entry [0][4]}, label={lst:c29-750e}]
#include <bits/stdc++.h>
using namespace std;
typedef long long ll;
const int K = 5;
const ll INF = (1LL<<60);
string P = "2017";
struct Mat { ll a[K][K];
    Mat(){ for(int i=0;i<K;i++) for(int j=0;j<K;j++) a[i][j]=INF; } };
Mat compose(const Mat&X, const Mat&Y){      // X then Y
    Mat Z;
    for(int i=0;i<K;i++) for(int k=0;k<K;k++) if(X.a[i][k]<INF)
        for(int j=0;j<K;j++) if(Y.a[k][j]<INF)
            Z.a[i][j]=min(Z.a[i][j], X.a[i][k]+Y.a[k][j]);
    return Z;
}
Mat make(char c){                            // one digit's transition matrix
    Mat M;
    for(int i=0;i<K;i++){
        M.a[i][i] = min(M.a[i][i], 1LL);     // delete this digit
        if(i<4 && c==P[i]) M.a[i][i+1] = min(M.a[i][i+1], 0LL);  // keep, advance
        else if(i>=3 && c=='6') { /* keep forbidden: forms 2016 */ }
        else M.a[i][i] = min(M.a[i][i], 0LL); // keep, inert
    }
    return M;
}
int n,q; char s[200005]; Mat tree_[800005];
void build(int i,int l,int r){
    if(l==r){ tree_[i]=make(s[l]); return; }
    int m=(l+r)>>1; build(i<<1,l,m); build(i<<1|1,m+1,r);
    tree_[i]=compose(tree_[i<<1], tree_[i<<1|1]);
}
Mat query(int i,int l,int r,int ql,int qr){
    if(ql<=l && r<=qr) return tree_[i];
    int m=(l+r)>>1;
    if(qr<=m) return query(i<<1,l,m,ql,qr);
    if(ql> m) return query(i<<1|1,m+1,r,ql,qr);
    return compose(query(i<<1,l,m,ql,qr), query(i<<1|1,m+1,r,ql,qr));
}
int main(){
    scanf("%d%d%s",&n,&q,s+1);
    build(1,1,n);
    while(q--){ int a,b; scanf("%d%d",&a,&b);
        Mat R = query(1,1,n,a,b);
        ll ans = R.a[0][4];                  // state 0 -> 4
        printf(ans>=INF ? "-1\n" : "%lld\n", ans);
    }
}
\end{lstlisting}

\begin{tracebox}[The greedy automaton on sample~\(\mathbf{1}\)'s first query]
Substring \(\texttt{20166766}\) (query \([1,8]\)).  Walk states, greedily keeping
whatever is inert or advancing and deleting only the dangerous \(\texttt{6}\)'s:
\texttt{2} advance \(0\!\to\!1\); \texttt{0} \(1\!\to\!2\); \texttt{1}
\(2\!\to\!3\); now \(\texttt{6},\texttt{6}\) are at state \(\ge3\Rightarrow\)
\textbf{delete both} (cost~\(2\)); \(\texttt{7}\) advances \(3\!\to\!4\); the last
two \(\texttt{6}\)'s at state~\(4\Rightarrow\) \textbf{delete} (cost~\(4\)).  Final
state~\(4\), deletions \(\boxed{4}\) -- exactly ``all four sixes must be removed''.
Query \([1,7]=\texttt{2016676}\): same until the trailing \(\texttt{6}\) is deleted
at state~\(4\) \(\Rightarrow 3\).  Query \([2,8]=\texttt{0166766}\) has no leading
\(\texttt{2}\): state never leaves~\(0\), so entry \([0][4]=\infty\Rightarrow
-1\).  The matrix product for \([1,8]\) reads \([0][4]=4\), matching \(\figref{c29-func-auto}\)'s
chain.
\end{tracebox}

\begin{warningbox}[Two order traps in one tree]
\((\min,+)\) product is \emph{not commutative}: \(\mathit{compose}(X,Y)\ne
\mathit{compose}(Y,X)\).  A node must store \(F_{\text{left}}\) then
\(F_{\text{right}}\), and a range fold must concatenate its \(O(\log n)\) cover
nodes strictly \emph{left to right} -- the split-node \code{query} above returns
\code{compose(left\_part, right\_part)}, never the reverse.  Second, the identity
matrix for ``empty'' ranges must be the \((\min,+)\) identity (\(0\) diagonal,
\(\infty\) elsewhere), \emph{not} all-zeros -- an all-zero matrix would let the
state jump for free and silently under-count deletions.
\end{warningbox}

\begin{hardeningbox}
\begin{itemize}[nosep,leftmargin=1.4em]
  \item \textbf{Reachability \(=\)-1, not a big number.} Test \(\mathit{ans}\ge
        \mathit{INF}\) before printing; \(\mathit{INF}\) plus a real cost must stay
        in \code{long long} (cap it or guard the inner \(\min\)).
  \item \textbf{The general \(K\).} For other required/forbidden subsequences
        \(K\) is the pattern length \(+1\); \(O(K^3)\) per compose means the method
        is for \emph{small} fixed patterns only.
  \item \textbf{Subsequence, not substring.} ``2017 as a subsequence'' is why the
        greedy longest-prefix state is enough; a contiguous-substring variant needs
        a KMP-style automaton (still a matrix, different states).
  \item \textbf{Same tree, affine variant.} Swap \code{Mat}/\code{compose} for
        \((a,b)\) maps to answer range affine-update-and-point-query -- the
        Ch.~20 \(\secref{c15-compose}\) tag reused as a stored aggregate.
\end{itemize}
\end{hardeningbox}

\begin{complexitybox}[Composition-tree cost]
Build \(O(n K^3)\) (a constant \(125\) here), point update \(O(K^3\log n)\), range
query \(O(K^3\log n)\) returning the composed map, evaluated on a state vector in
\(O(K^2)\).  For CF~750E, \(K=5\): \(O((n+q)\cdot125\cdot\log n)\), comfortably
inside the limits.  The whole cost is the \(K^3\) matrix product at each of the
\(O(\log n)\) levels.
\end{complexitybox}

\begin{recallbox}[23.4 in one breath]
\emph{A composition tree stores a \emph{map} at each node and merges by
function composition (\(\defref{c29-func}\), \(\thmref{c29-func-mono}\)) --
Ch.~20's lazy monoid promoted to the aggregate.  Affine maps \((a,b)\) compose to
\((a_1a_2,\,a_1b_2+b_1)\); a finite-state machine composes by \((\min,+)\) matrix
product (\(\coderef{c29-func}\)).  CF~750E's ``\(2017\) present, \(2016\) absent''
is a \(5\)-state automaton (\(\figref{c29-func-auto}\)) -- delete the dangerous
\(\texttt{6}\)'s, advance on the pattern symbols -- so a query is one ordered
matrix product read at \([0][4]\).}
\selftest{\item What algebraic law does a composition tree need of its merge, and
what does it \emph{not} need? \emph{(Associativity; it does not need commutativity
-- order of composition is semantically real.)}
\item Why is a kept \(\texttt{6}\) charged as a delete in 750E exactly at state
\(\ge3\)? \emph{(State \(\ge3\) means ``201'' is already a kept subsequence, so a
\(\texttt{6}\) would create the forbidden \(\texttt{2016}\).)}}
\end{recallbox}
